\subsection{Requisitos No Funcionales}

Los requisitos no funcionales (RNF) establecen propiedades de calidad y restricciones transversales del producto, con independencia de la tecnología usada para implementarlas.

\subsubsection*{\textbf{RNF-01} - Acceso restringido y diferenciación de capacidades}

El acceso debe limitarse según las atribuciones vigentes, el ámbito institucional aplicable y, para la consulta clínica del personal asistencial, la relación asistencial con la hospitalización. Las capacidades disponibles deben diferenciarse según la autorización efectiva de cada persona.

\subsubsection*{\textbf{RNF-02} - Confidencialidad y privacidad}

La información debe protegerse durante su transmisión, almacenamiento y procesamiento. El procesamiento de capacidades asistidas no debe enviar contenido clínico a proveedores externos de inteligencia artificial.

\subsubsection*{\textbf{RNF-03} - Integridad de la información}

El sistema debe preservar la procedencia de la información, distinguir el contenido proveniente de fuentes externas del contenido propio o generado y prevenir duplicados y modificaciones no autorizadas.

\subsubsection*{\textbf{RNF-04} - Trazabilidad}

Las operaciones y el contenido registrado deben conservar metadatos suficientes para identificar procedencia, autoría, tiempos relevantes y operaciones realizadas, según corresponda.

\subsubsection*{\textbf{RNF-05} - Seguridad del contenido asistido}

El contenido asistido no debe incorporar antecedentes no presentes en sus fuentes ni formular diagnósticos. No debe publicarse automáticamente y debe quedar sujeto a revisión humana antes de cualquier uso como nota.

\subsubsection*{\textbf{RNF-06} - Usabilidad y baja carga cognitiva}

La interfaz móvil debe facilitar las tareas principales del personal asistencial con una presentación clara y una carga cognitiva reducida, sin imponer un modo de interacción no sustentado por el contexto de uso.

\subsubsection*{\textbf{RNF-07} - Rendimiento operativo}

El rendimiento de las operaciones relevantes debe ser compatible con las tareas previstas y sus criterios de aceptación. Debe evaluarse mediante mediciones reproducibles en dispositivos representativos del entorno objetivo e informarse con sus condiciones de prueba, sin presuponer umbrales no medidos.

\subsubsection*{\textbf{RNF-08} - Interoperabilidad}

Las interfaces de intercambio deben definir contratos explícitos para las fuentes integradas y permitir mapear la información pertinente al modelo del sistema. La selección de estándares debe responder al contrato de cada integración y no presupone un estándar único obligatorio.

\subsubsection*{\textbf{RNF-09} - Manejo de fallos}

Los fallos de una operación deben preservar las fuentes de información y no impedir tareas independientes que puedan continuar de forma segura. La recuperación ante fallos nunca debe omitir controles de autorización ni revisión humana.

\subsubsection*{\textbf{RNF-10} - Retención mínima}

Los artefactos temporales deben conservarse solo durante el tiempo necesario para completar su propósito y eliminarse conforme a la política de retención aplicable. La política debe contemplar los residuos de operaciones interrumpidas.
